\subsection{Workload Methodology and Coverage}
\label{sec:workload-methodology}

We evaluate a four-core, 1-GHz NPU in which each core contains a
$128\times128$ weight-stationary systolic array.  Table~\ref{tab:workload-suite}
summarizes the evaluated workload configurations.  All image-classification
graphs use batch size one and execute end to end.  GPT-2 also executes its
complete optimized ONNX graph with batch size one, sequence length 1024,
past-sequence length zero, and output-sequence length 1025.

The language-model generator is evaluated with five requests.  Each request
has prompt length 32, target length 8, and cached length zero; request arrivals
occur at 0, 100,000, 200,000, 300,000, and 400,000 cycles.  The scheduler is
configured with a maximum batch size of eight.  We refer to this experiment as
the \emph{32/8 configuration}.  OPT-125M is evaluated both as a complete
12-layer model and as one representative decoder layer.  The larger OPT and
LLaMA models use one representative decoder layer because their complete-model
execution cost scales with layer count.  The representative-layer claim is
validated separately below and is restricted to the repeated OPT decoder
structure.

\begin{table*}[t]
  \centering
  \caption{Evaluated workload suite and execution scope.  ``Full'' denotes an
  end-to-end graph or all model layers; ``1/$L$'' denotes one representative
  layer from an $L$-layer language model.}
  \label{tab:workload-suite}
  \begin{tabular}{lllcl}
    \toprule
    Workload(s) & Family & Input configuration & Scope & Model path \\
    \midrule
    AlexNet; ResNet-18/34/50/101 & CNN & batch 1 & Full & Audited ONNX \\
    VGG-11/16/19; Wide-ResNet-50-2 & CNN & batch 1 & Full & Audited ONNX \\
    GPT-2 & Transformer & batch 1, sequence 1024 & Full & Audited ONNX \\
    OPT-125M & Transformer & five 32/8 requests, max batch 8 & Full (12 layers) & Native generator \\
    OPT-1.3B & Transformer & five 32/8 requests, max batch 8 & 1/24 & Native generator \\
    OPT-2.7B & Transformer & five 32/8 requests, max batch 8 & 1/32 & Native generator \\
    OPT-6.7B & Transformer & five 32/8 requests, max batch 8 & 1/32 & Native generator \\
    OPT-66B & Transformer & five 32/8 requests, max batch 8 & 1/64 & Native generator \\
    LLaMA-3-8B & Transformer & five 32/8 requests, max batch 8 & 1/32 & Native generator \\
    \bottomrule
  \end{tabular}
\end{table*}

For imported ONNX graphs, we perform a strict operator audit before simulation.
The modeled set comprises convolution, fused convolution, matrix multiplication
and GEMM variants, pooling, flattening, attention, embedding normalization,
skip-layer normalization, and fused GELU variants.  Cast and output-reorder
nodes are accepted only as explicit pass-through operations.  A graph is
rejected if it contains any other operator, because ONNXim would otherwise
represent that node as a zero-work dummy operation.  Grouped or depthwise
convolution is also rejected by the current weight-stationary mapping path.
An audit pass is necessary but not sufficient: the run must additionally
initialize, complete, satisfy mapping capacity, and pass the simulator's final
accounting checks.  The ten graph workloads in Table~\ref{tab:workload-suite}
satisfied these requirements.  Language models generated directly by ONNXim
do not pass through the ONNX operator-audit path.

\subsubsection{Representative-layer validation}
\label{sec:representative-layer-validation}

Table~\ref{tab:representative-layer} compares one OPT-125M decoder layer with
the complete 12-layer execution under the same 32/8 configuration.  The
complete/single-layer ratios closely track the ideal factor of 12.  In
particular, the normalized net saving changes from 1.1182\% to 1.1220\%, a
difference of only 0.0038 percentage points.  This result supports
representative-layer characterization for structurally repeated OPT decoder
blocks under this configuration.  It does not establish the same equivalence
for unrelated architectures or input configurations.  Only the full
OPT-125M result is included in the main workload aggregate.

\begin{table}[t]
  \centering
  \footnotesize
  \setlength{\tabcolsep}{3pt}
  \caption{Single-layer versus complete 12-layer OPT-125M validation.}
  \label{tab:representative-layer}
  \begin{tabular}{lrrr}
    \toprule
    Metric & One layer & 12 layers & Full/single \\
    \midrule
    Core-cycles & 4,196,744 & 50,194,496 & 11.9603 \\
    OFF cycles & 515,516 & 6,210,937 & 12.0480 \\
    OFF coverage (\%) & 12.2837 & 12.3737 & -- \\
    Gross saving (J) & $4.52584\!\times\!10^{-7}$ & $5.42518\!\times\!10^{-6}$ & 11.9871 \\
    Controller energy (J) & $1.82071\!\times\!10^{-7}$ & $2.17879\!\times\!10^{-6}$ & 11.9667 \\
    Net saving (J) & $2.70513\!\times\!10^{-7}$ & $3.24640\!\times\!10^{-6}$ & 12.0009 \\
    Net/reference (\%) & 1.1182 & 1.1220 & -- \\
    \bottomrule
  \end{tabular}
\end{table}

\subsection{Workload-Level Energy Results}
\label{sec:workload-energy-results}

For each workload, exact synthesized-controller replay supplies the OFF
episodes.  The direct HSPICE duration--energy curve supplies the corresponding
array transition and leakage energy, while 100\%-annotated ASAP7 gate-level
activity supplies controller energy.  Let $P_{\mathrm{on}}$ be the measured
steady-state ON leakage of the scaled $128\times128$ array and $T_w$ the sum of
core-cycles for workload $w$.  The common normalization reference is
\begin{equation}
  E_{\mathrm{ref},w}=P_{\mathrm{on}}\frac{T_w}{10^9},
  \qquad P_{\mathrm{on}}=5.7643057152\ \mathrm{mW}.
\end{equation}
Gross saving includes array leakage reduction minus the measured OFF/ON event
cost.  Net saving subtracts the full-trace controller energy:
\begin{equation}
  E_{\mathrm{net},w}=E_{\mathrm{gross},w}-E_{\mathrm{ctrl},w}.
\end{equation}
These percentages therefore quantify saving relative to the array's
always-ON static-energy reference; they are not whole-chip energy reductions.

\begin{table*}[t]
  \centering
  \caption{Measured workload results.  OFF is the fraction of total
  core-cycles switched OFF.  Energy columns are percentages of the always-ON
  SA static-energy reference.  OPT-66B net energy is omitted because only one
  of four required controller-power reports is retained.}
  \label{tab:workload-energy}
  \begin{tabular}{lrrrr}
    \toprule
    Workload & OFF cycles (\%) & Gross saving (\%) & Controller (\%) & Net saving (\%) \\
    \midrule
    AlexNet             & 11.773 & 1.317 & 0.688 & +0.630 \\
    ResNet-18           &  0.000 & 0.000 & 0.576 & -0.576 \\
    ResNet-34           &  0.000 & 0.000 & 0.577 & -0.577 \\
    ResNet-50           &  0.000 & 0.000 & 0.591 & -0.591 \\
    ResNet-101          &  0.000 & 0.000 & 0.593 & -0.593 \\
    VGG-11              & 17.348 & 0.960 & 0.752 & +0.209 \\
    VGG-16              & 20.824 & 1.398 & 0.735 & +0.663 \\
    VGG-19              & 17.840 & 1.125 & 0.715 & +0.410 \\
    Wide-ResNet-50-2    &  0.000 & 0.000 & 0.592 & -0.592 \\
    GPT-2               &  1.342 & 0.205 & 0.681 & -0.475 \\
    OPT-125M (12 layers)& 12.374 & 1.875 & 0.753 & +1.122 \\
    OPT-1.3B            & 26.582 & 1.835 & 0.813 & +1.022 \\
    OPT-2.7B            & 28.491 & 1.840 & 0.831 & +1.009 \\
    OPT-6.7B            & 24.702 & 1.745 & 0.821 & +0.924 \\
    OPT-66B             & 22.386 & 2.000 & --    & -- \\
    LLaMA-3-8B          & 23.635 & 1.638 & 0.819 & +0.819 \\
    \bottomrule
  \end{tabular}
\end{table*}

Among the 15 workloads with complete net-energy evidence, nine are positive
and six are negative.  Their equal-workload arithmetic mean is +0.227\% and
their median is +0.410\%.  Weighting instead by the static-reference energy
(equivalently, by total core-cycles at the fixed array leakage and frequency)
gives +0.723\%.  The two aggregates answer different questions and are both
reported explicitly to avoid allowing long workloads to dominate an
unqualified average.

The signs follow directly from the measured balance between opportunity and
overhead.  ResNet-18/34/50/101 and Wide-ResNet-50-2 have no certified OFF
cycles, so their gross saving is zero while the controller remains active.
GPT-2 exposes only 1.342\% OFF coverage: its 0.205\% gross reduction is smaller
than its 0.681\% controller cost.  Conversely, the VGG, OPT, LLaMA, and AlexNet
cases provide sufficient profitable OFF intervals to exceed controller
energy.  Zero certified OFF coverage does not imply that the SA has no raw
idle cycles; it means that those idle cycles do not satisfy the present
causal, break-even, and non-retentive safety conditions.

\subsection{Limitations and Proposed Retention Extension}
\label{sec:limitations-retention}

The current evaluation has the following limitations.
\begin{itemize}
  \item Power-state decisions are replayed as a shadow mechanism and do not
  feed wake-up delay back into ONNXim scheduling.  The current results
  therefore establish cycle-preserving decisions for the evaluated traces,
  but do not measure a performance--energy trade-off with nonzero wake delay.
  \item The mechanism is non-retentive.  It permits OFF only when the next
  certified SA use reloads state through \texttt{GEMM\_PRELOAD}; intervals
  leading directly to a state-reusing \texttt{GEMM} are excluded.
  \item The M300-equivalent header is a characterized design point, not a
  proven full-array optimum.  The model does not include a routed full-array
  power-delivery network or worst-case instantaneous droop.
  \item Array event energy is scaled linearly per PE from direct $1\times1$
  characterization, with a $4\times4$ experiment used as the available
  scaling check.  A transistor-level $128\times128$ array is not simulated.
  \item Controller power is synthesis-stage, 100\%-annotated gate power.  The
  reported implementation excludes placed power-grid parasitics, isolation
  cells, routed clock-tree and interconnect power, delay-dependent glitch
  power, and architectural reset/control energy.  The descriptor-generation
  hardware that would replace simulator-side metadata production is also not
  included in the synthesized controller area or power.
  \item Results cover one four-core, $128\times128$, 1-GHz NPU configuration.
  Language inputs use the stated 32/8 request stream, graph inputs use batch
  one, and unsupported ONNX graphs are excluded.  Representative-layer
  equivalence is measured for OPT-125M only and is not assumed for all model
  families.
  \item OPT-66B has complete replay and gross array-energy evidence, but lacks
  three of the four controller-power reports required for an exact net result;
  it is excluded from all net-energy aggregates.
\end{itemize}

A state-retentive extension is a proposed next step, not a measured result.
It would preserve all PE state required by a subsequent state-reusing
\texttt{GEMM}, allowing the controller to consider intervals that the present
non-retentive certificate rejects.  The existing causal lower bounds and
frozen countdown can remain, while the policy selects among ON, retention,
and OFF states: retention for intervals requiring preserved state, and OFF
for sufficiently long intervals whose next use performs a certified preload.

Quantifying this extension requires new circuit and architectural evidence:
retention voltage and leakage, entry and wake energy, entry and wake latency,
and state-correctness characterization; synthesis and activity-based power for
the additional mode-control logic; simulator feedback of wake delay to measure
performance loss; and replay or formal validation of every state transition.
Until those measurements are available, no retention energy saving or
performance penalty is claimed.
